% =====================================================================
\chapter{Planilha de apoio}\label{an:planilha}
% =====================================================================

A planilha de apoio é entregue em arquivo eletrônico, com o nome
\texttt{TCC Atual - Cenário 3 (corrigido).xlsx}, acompanhada das versões dos
Cenários 1 e 2, usadas nos testes de robustez. Ela contém as seguintes abas:

\begin{alineas}
\item \textbf{Saúde - 1218 à 1225:} base da Economática com as 16 companhias e os
indicadores trimestrais de 31/12/2018 a 31/12/2025;
\item \textbf{Classificação, Estudos, Resumo dos Critérios, Metodologia e Cenários:}
critérios de classificação e sua aplicação a cada empresa;
\item \textbf{Indicadores e Metodologias:} definição dos indicadores e explicação dos
testes estatísticos;
\item \textbf{Análises:} médias anuais por grupo, testes do período completo e testes
por subperíodo, para o indicador escolhido na célula B1;
\item \textbf{uma aba por indicador:} médias trimestrais por grupo, valores por
empresa e testes sobre a média anual de cada empresa.
\end{alineas}

% =====================================================================
\chapter{Código dos cálculos}\label{an:script}
% =====================================================================

Os resultados, as tabelas e as figuras do \autoref{cap:resultados} foram gerados
pelos programas a seguir, escritos em Python, que leem a aba-base da planilha do
\autoref{an:planilha}. Os programas foram elaborados com o auxílio de ferramenta
de inteligência artificial generativa (Claude, da Anthropic); o autor definiu os
procedimentos e conferiu os resultados.

\section*{resultados\_c3.py}
\lstinputlisting{anexos/resultados_c3.py}

\section*{robustez\_extra.py}
\lstinputlisting{anexos/robustez_extra.py}

\section*{ajustes\_banca.py}
\lstinputlisting{anexos/ajustes_banca.py}

\section*{graficos.py}
\lstinputlisting{anexos/graficos.py}
